% !TeX root = ../main.tex
\chapter{Conclusion and Outlook}
\label{ch:conclusion}

This chapter closes the thesis by returning to the objective of
\cref{sec:objectives}: integrate whole-body \mpc{} on the physical
\hero{} platform and characterise its \ee{}-tracking behaviour. The
preceding chapters have already separated what the optimiser does on its
kinematic model (\cref{ch:simulation}), what the MRT bridge adds
(\cref{sec:hw-challenges}), and what the plant adds on top
(\cref{sec:hw-results}). What follows is a concise account of what was
achieved, what the present data do and do not establish, and which
engineering steps would most directly change the numbers reported in
\cref{ch:hardware}.

% ---------------------------------------------------------------------
\section{Summary}
\label{sec:concl-summary}

The work set out to move \ocs{} whole-body \mpc{} from a framework
demonstration to a deployable stack on a differential-drive mobile
manipulator in an LMAS-oriented setting. Against that goal, four
contributions were delivered.

First, a whole-body motion-planning approach was selected for the
\hero{} platform after comparing \ocs{} with alternative frameworks
(\cref{ch:sota}). Second, \ocs{} was integrated into the existing
software stack: a split host/on-board deployment over Zenoh keeps the
actuation path local to the IPC while the SLQ solver runs on a
workstation (\cref{sec:distributed-arch}); a nine-state kinematic model
with non-holonomic base, safety-panel joint limits, and a reduced
self-collision set was parameterised for on-robot operation (\cref{sec:wholebody-model});
and an MRT bridge turns the published policy into wheel and joint
commands under explicit guards (\cref{sec:hw-interfaces}). Third, a
simulation-based acceptance path was built around Experiment~A: the same
task file and waypoint list that the robot will load are exercised in a
dummy MRT loop, with internal metrology and ISO-oriented path metrics
defined before any hardware session (\cref{ch:simulation}). Fourth, that
configuration was taken through mock \texttt{ros2\_control} hardware and
onto the physical platform. The closed eight-corner path of
\cref{fig:sim-box} completed on both plants; the approaching gate,
segment tagging, and lead-out hold that \cref{sec:sim-node} specifies all
appeared in the recorded sessions.

The tracking outcome is more qualified. On mock hardware, which shares
the bridge and the policy lag with the robot but replaces the plant by
command echoing, the attained \ee{} path followed the commanded rectangle
at a few centimetres (\cref{tab:hw-metrics}). On the physical robot, with
the same files and the same wall-clock lap duration, position error
remained at one to two decimetres on the long edges and did not collapse
during corner dwells (\cref{fig:hw-error,tab:hw-segments}). The loop
therefore ran in real time in the sense of \cref{eq:rt-invariant}---no
policy-exhaustion hold occurred on the hardware lap---but the effective
\mpc{} rate was about \qty{4.4}{\hz} against a \qty{10}{\hz} request, and
the first one to two tenths of a second of each plan were lost to solve
time and network delay before the bridge could apply them
(\cref{fig:hw-lag}). Along-track lag consistent with that delay explains
the mock residual; it does not explain the larger hardware offset
(\cref{sec:hw-challenges}).

Several constraints bound how strongly these results can be read. All
desired and attained poses were reconstructed from wheel odometry and
joint encoders inside the waypoint node (\cref{sec:sim-metrology}), so
the scores measure internal consistency of the closed loop, not absolute
accuracy in the hall. Each plant was exercised once, at
$v_{\mathrm{ref}}=\qty{0.05}{\meter\per\second}$, without the three
mandatory ISO speed levels or repeated cycles. The workstation was
connected by Ethernet throughout (\cref{sec:hw-deployment}); the numbers
therefore do not speak to control over the platform's wireless link. The
\ocs{} model treats the base as a velocity integrator and knows nothing
of the drive controller's acceleration limit or of the UR10e following
error under streaming position control. Finally, no decoupled Nav2 plus
MoveIt baseline was run on the same path, so the thesis validates
whole-body execution and characterises its residual; it does not yet
demonstrate superiority over sequential planning on identical metrics.

% ---------------------------------------------------------------------
\section{Critical Reflection}
\label{sec:concl-reflection}

The experiments answer a narrower question than the LMAS motivation of
\cref{sec:lmas} might suggest. They show that a time-parameterised
\ee{} reference, published as \ocs{} \texttt{TargetTrajectories} rather
than as a sequence of static goals, can be executed end to end on
\hero{} hardware by a receding-horizon whole-body controller, and they
localise the largest tracking gap to the plant and to model mismatch
rather than to a broken data pipeline or a reference that the dummy
already rejects. They do not, by themselves, establish that this
pipeline meets production accuracy requirements, nor that whole-body
\mpc{} outperforms a well-tuned decoupled stack on task completion time
or corner accuracy.

The metric layer deserves the same caution. Time-aligned pose error
$e_{\mathrm{p}}$ and the along-track lag $t_{\mathrm{lag}}$ of
\cref{eq:lag-eff} are meaningful comparisons \emph{within} one
controller, because desired and attained poses share the same
observation time and the same state estimate. Hausdorff distance and its
mean and percentile variants (\cref{eq:mean-hausdorff}) remove the
assumption of a common speed profile and are the right preparatory step for a later
comparison with a pose-to-pose baseline, but with $n=1$ cycle they still
do not yield ISO path accuracy $\mathrm{AT}_{p}$ or repeatability
$\mathrm{RT}_{p}$ in the sense of ISO~9283~\cite{ISO9283}. An external
measurement of \texttt{tool0}---for example from an OptiTrack rig as
already noted in \cref{sec:sim-metrology}---would be needed before any
claim about absolute \ee{} pose error enters the thesis.

The split between mock and hardware results is instructive in what it
rules out. Mock hardware preserves policy lag, velocity clamps, and the
Zenoh hop; when it tracks well, the reference geometry, the \ocs{} time
axis, and the cost tuning are not the dominant fault. Hardware error that
grows on horizontal edges where the base must contribute, persists when
the command stops at a dwell, and exceeds a pure lag prediction of
$0.36\,\mathrm{s}\times v_{\mathrm{ref}}$ (\cref{sec:hw-results}),
points instead to unmodelled base dynamics and to re-solves that start
from a state the previous plan did not anticipate. That reading is
consistent with the kinematic integrator in the model and with the
\qty{400}{\kilo\gram} platform's own motion limits, but it remains a
diagnosis from internal signals; isolating wheel slip, odometry drift,
and arm tracking error would need additional sensing or deliberately
simpler motions.

Scope limits are finally worth stating plainly. One mobile manipulator,
one closed path, and one lap per plant are enough to debug integration
and to expose the main failure modes; they are not enough to support
statistical claims about repeatability or about behaviour across the
workspace. The progression from RViz single-goal regulation to
time-parameterised trajectory tracking via the waypoint node
(\cref{sec:sim-node}) is a real extension of the upstream \ocs{} demo,
but the evaluation has so far exercised only the tracking regime on a
single box geometry.

% ---------------------------------------------------------------------
\section{Recommendations for Resolving Present Issues}
\label{sec:concl-recommendations}

The hardware lap already identifies where effort would change the
recorded numbers, rather than merely widen the solution time window.

\paragraph{Raise the effective \mpc{} rate and shorten policy lag.}
Mean solve time on the hardware session was \qty{227}{\milli\second}
(\cref{tab:hw-metrics}), so the controller replans at roughly half the
configured \qty{10}{\hz}. The DDP benchmark that \ocs{} prints at
shutdown attributes that time to two stages above all. In an earlier
instrumented session on the same node, summarised over $1794$
iterations, \emph{Initialization} averaged \qty{170}{\milli\second} per
iteration (\qty{56}{\percent} of solver runtime) and \emph{Search
Strategy} \qty{111}{\milli\second} (\qty{36}{\percent}), while the
backward pass stayed below \qty{1}{\milli\second}. The absolute values
belong to that session rather than to the laps of \cref{ch:hardware},
but the split suggests attacking warm-start cost and rollout search
before adding SLQ iterations. Pinning the solver threads
away from unrelated workstation load, keeping the Release build used here,
and revisiting the twelve-pair self-collision set if distance queries
still dominate are practical first steps. Enabling a feedback policy
(\texttt{useFeedbackPolicy=true} in the task file; currently disabled in
\cref{tab:mpc-config}) would let the bridge use linear gains between
sparse policy updates; the trade-off is larger policy messages over
Zenoh, which must be measured on the link that will actually be used in
production. Widening $T_{\mathrm{sol}}$ further, as in the Experiment~A
file, prevents exhaustion when a solve spikes; it does not restore the
aggressive $\tau=0$ input that lag already consumes (\cref{fig:hw-lag}).

\paragraph{Close the model--plant gap on the base.}
The largest hardware residual appears where the base must translate while
the arm holds orientation, and it does not relax when the reference stops.
Adding an acceleration limit or a first-order lag on base velocities
inside the \ocs{} model would align the optimiser with what the
freeMoVe drive stack enforces and reduce the state mismatch that each
re-solve inherits. A complementary experiment---base motion on a closed
planar path with the \ee{} held fixed in space, the converse of
Experiment~A---would separate base tracking from arm-dominated edges
without changing the software architecture.

\paragraph{Improve pose estimation and external metrology.}
If odometry drift dominates closed-loop drift, fusing wheel odometry
with additional proprioceptive or exteroceptive sources through an EKF
(for example \texttt{robot\_localization}) is the standard next step; the
present pipeline uses raw wheel odometry only (\cref{sec:hero-platform}).
For claims that must survive examination against ISO~9283, an independent
\ee{} pose remains mandatory. Ground-truth capture would also break the
confounding between desired and attained frames noted in
\cref{sec:sim-metrology}.

\paragraph{Complete the evaluation programme.}
Three items would turn the present characterisation into a comparative
study aligned with \cref{sec:objectives}. A decoupled Nav2 plus MoveIt
baseline on the same corner set would supply $\mathrm{AP}_{p}$ and task
time on common footing (\cref{sec:sim-path-metrics}). ISO Table~3 speed
levels and $n\ge 3$ cycles on the whole-body pipeline would populate
$\mathrm{AT}_{p}$ and $\mathrm{RT}_{p}$. Re-running the mock versus
hardware comparison over a wireless Zenoh link would test whether the
Ethernet sessions of \cref{sec:hw-deployment} bias the timing budget
optimistically. Until those pieces exist, the defensible statement is the
one \cref{sec:hw-results} already makes: the Experiment~A files execute
on hardware; mock hardware tracks at centimetre scale with lag consistent
with policy delay; the physical platform does not, and the gap is not
explained by solver failure alone.
